\documentclass[11pt]{article}
\usepackage[utf8]{inputenc}
\usepackage{amsmath,amssymb,amsthm}
\usepackage[margin=1in]{geometry}
\usepackage{hyperref}
\usepackage{xcolor}
\usepackage{enumitem}
\usepackage{newunicodechar}
\newunicodechar{ℝ}{\ensuremath{\mathbb{R}}}
\newunicodechar{⟨}{\ensuremath{\langle}}
\newunicodechar{⟩}{\ensuremath{\rangle}}
\newunicodechar{α}{\ensuremath{\alpha}}
\newunicodechar{π}{\ensuremath{\pi}}
\newunicodechar{Γ}{\ensuremath{\Gamma}}
\newunicodechar{→}{\ensuremath{\to}}
\newunicodechar{ℕ}{\ensuremath{\mathbb{N}}}
\newunicodechar{₁}{\ensuremath{_1}}
\newunicodechar{₂}{\ensuremath{_2}}
\newunicodechar{𝓝}{\ensuremath{\mathcal{N}}}
\newunicodechar{≤}{\ensuremath{\leq}}
\newunicodechar{∀}{\ensuremath{\forall}}
\newunicodechar{∫}{\ensuremath{\int}}
\newunicodechar{∞}{\ensuremath{\infty}}
\newunicodechar{≠}{\ensuremath{\neq}}

\title{Tide retrospective:\\J-function asymptotic-equivalence wrapper}
\author{Daniel Murfet}
\date{7 May 2026}

\begin{document}
\maketitle

\begin{abstract}
This is a 200-line tide that closes the J-function programme's
narrative: with the parametric form $t^{1/(2k)} \cdot J_k(t) \to \pi(0)
\cdot \mathrm{CONST}_k$ already established by the generic-$k$ tide
(N1) and the matching partition $Z_k(t)$ available in closed form,
the asymptotic-equivalence statement actually used in the grammar
paper,
\[
  J_k(t) \;\sim\; \pi(0) \cdot Z_k(t) \quad (t \to \infty),
\]
collapses to a thin \texttt{Tendsto.div\_const} corollary plus a
Mathlib idiom-bridge.\footnote{Specifically
\texttt{Asymptotics.isEquivalent\_iff\_tendsto\_one}.} The
unconditional ratio form $J_k/Z_k \to \pi(0)$ is the cleaner
statement and the building block; the \emph{IsEquivalent} form is the
named user-facing corollary under $\pi(0) \neq 0$. Single session,
GPT-5.5~Pro and Claude voted Candidate~B (both theorems) on the first
round, zero-friction proof, 0 \texttt{sorry}.
\end{abstract}

\section{Setting}

The grammar-paper precursor sequence has three landings on 7 May:
quartic ($k=2$, the E4 tide), sextic ($k=3$, the L1 tide), and
generic-$k$ (the N1 tide). Each established the same shape of
asymptotic in a $t^{1/(2k)}$-rescaled form:
\[
  t^{1/(2k)} \cdot J_k(t) \;\longrightarrow\; \pi(0) \cdot
  \tfrac{1}{k} \cdot ((2k)!)^{1/(2k)} \cdot \Gamma(1/(2k)),
\]
which is mathematically equivalent to (but typographically distinct
from) the form the grammar paper actually uses:
\[
  J_k(t) \;\sim\; \pi(0) \cdot Z_k(t),
\]
where $Z_k(t) = \int e^{-tw^{2k}/(2k)!}\,dw$ is the matching
partition. The L1 tide explicitly named this conversion as a
follow-up: \emph{``The grammar paper's natural statement is $J(t) \sim
\pi(0) \cdot Z(t)$ \ldots\ Conversion is \texttt{Tendsto.div} against
\texttt{quartic\_partition}'s positive-and-tending-to-zero $Z(t)$.
$\sim$10 lines on top of \texttt{JFunctionQuartic.lean}.''}

The L1 estimate of $\sim$10 lines was for the $k=2$ specialisation.
After N1, the right place to land the wrapper is at the
\emph{parametric} level — once on
\texttt{JFunctionMonomial.lean}, valid for all $k \geq 1$, with the
$k=2$ and $k=3$ specialisations as zero-line corollaries (instantiate
$k$). The actual implementation took $\sim$200 lines, an overrun
driven by closed-form-positivity bookkeeping rather than by any
mathematical surprise.

This tide is therefore another \emph{narrative-closer}: it does not
prove a new asymptotic content; it packages the existing asymptotic
into the named idiom the paper uses.

\section{The thing formalised}

\begin{quote}\itshape
For the single-monomial potential $K(w) = w^{2k}/(2k)!$ with $k \geq
1$ and a continuous, globally bounded prior $\pi$, the J-function
\[
  J_k(t) := \int_{\mathbb{R}} e^{-t\,w^{2k}/(2k)!}\,\pi(w)\,dw
\]
and the matching partition
\[
  Z_k(t) := \int_{\mathbb{R}} e^{-t\,w^{2k}/(2k)!}\,dw
\]
satisfy two convergence statements as $t \to \infty$:
\begin{enumerate}[leftmargin=*]
\item (Unconditional ratio.) $\displaystyle \frac{J_k(t)}{Z_k(t)} \;\longrightarrow\; \pi(0)$.
\item (Asymptotic-equivalence.) Under $\pi(0) \neq 0$,
$J_k \sim \pi(0) \cdot Z_k$ in the sense that $J_k - \pi(0)\,Z_k = o(\pi(0)\,Z_k)$.
\end{enumerate}
\end{quote}

The Lean signatures, in
\texttt{Laplace/OneD/JFunctionAsymptoticEquivalence.lean}\footnote{The
file extends \texttt{JFunctionMonomial.lean} (the N1 seabed); it
imports \texttt{Mathlib.Analysis.Asymptotics.AsymptoticEquivalent} for
the named idiom \texttt{IsEquivalent}. Permalink: the file at the new
commit on the tide branch.}:

\begin{verbatim}
theorem kth_jfunction_div_partition
    {k : ℕ} (hk : 1 ≤ k)
    {prior : ℝ → ℝ} (hprior_cont : Continuous prior)
    {M : ℝ} (hprior_bd : ∀ x, |prior x| ≤ M) :
    Tendsto (fun t : ℝ =>
        (∫ w, exp(-(t * w^(2k) / (2k)!)) * prior w) /
        (∫ w, exp(-(t * w^(2k) / (2k)!))))
      atTop (𝓝 (prior 0))

theorem kth_jfunction_isEquivalent
    {k : ℕ} (hk : 1 ≤ k)
    {prior : ℝ → ℝ} (hprior_cont : Continuous prior)
    {M : ℝ} (hprior_bd : ∀ x, |prior x| ≤ M)
    (hprior0 : prior 0 ≠ 0) :
    (fun t : ℝ => ∫ w, exp(-(t * w^(2k) / (2k)!)) * prior w)
      ~[atTop]
    (fun t : ℝ => prior 0 * ∫ w, exp(-(t * w^(2k) / (2k)!)))
\end{verbatim}

The first is the building block; the second is the user-facing form.

\section{Proof strategy}

The core observation is short: the partition $Z_k(t)$ has an
\emph{exact} closed form, not just an asymptotic. From N1's
\texttt{kth\_partition} we get
\[
  Z_k(t) \;=\; \frac{1}{k}\,((2k)!/t)^{1/(2k)}\,\Gamma(1/(2k))
  \quad\text{for } t > 0,
\]
which, with $\alpha := 1/(2k)$ and $\mathrm{CONST}_k := (1/k) \cdot
((2k)!)^\alpha \cdot \Gamma(\alpha)$, rearranges to
\[
  t^\alpha \cdot Z_k(t) \;=\; \mathrm{CONST}_k \qquad \text{(exact, for all } t > 0\text{)}.
\]
Combined with N1's \texttt{kth\_jfunction\_asymptotic}, which gives
$t^\alpha \cdot J_k(t) \to \pi(0) \cdot \mathrm{CONST}_k$, the ratio is
\[
  \frac{J_k(t)}{Z_k(t)} \;=\; \frac{t^\alpha \cdot J_k(t)}{t^\alpha \cdot Z_k(t)}
  \;=\; \frac{t^\alpha \cdot J_k(t)}{\mathrm{CONST}_k}
  \;\longrightarrow\; \frac{\pi(0) \cdot \mathrm{CONST}_k}{\mathrm{CONST}_k}
  \;=\; \pi(0).
\]
Implementation in Lean: \texttt{Tendsto.div\_const} on the existing
N1 limit, plus a \texttt{Tendsto.congr'} to swap the form.

The \emph{IsEquivalent} corollary then uses Mathlib's
\texttt{isEquivalent\_iff\_tendsto\_one}: under eventual non-vanishing of
the target $\pi(0) \cdot Z_k(t)$, $f \sim g$ is equivalent to $f/g \to
1$. The eventual-nonvanishing condition reduces to $\pi(0) \neq 0$
(since $Z_k(t) > 0$ for all $t > 0$ by the closed form, which is a
positive product of positive factors). The ratio convergence
$f/g \to 1$ is then $(J_k/Z_k)/\pi(0) \to \pi(0)/\pi(0) = 1$, which
follows from the unconditional ratio theorem by a single
\texttt{Tendsto.div\_const}, modulo a \texttt{Pi.div} reshape.

The mathematical narrative is unsurprising; the Lean encoding is
mechanical. The 200-line file size comes from the closed-form
positivity bookkeeping (positivity of the constant, of the rpow, of
the Gamma function, of $(2k)!/t$ for $t > 0$) and from the
constant-massaging \texttt{set} block at the top of the unconditional
theorem, not from any clever proof.

\section{Roadblocks and resolutions}

Three small Lean-level frictions, all anticipated by the laplace
\texttt{CLAUDE.md}.

\paragraph{One transient ``no goals'' from over-eager
\texttt{field\_simp}.} The closed-form-rearrangement
$t^\alpha \cdot Z_k(t) = \mathrm{CONST}_k$ chains
\texttt{rw [kth\_partition hk ht]} with
\texttt{rw [Real.div\_rpow hfacnn htnn]} and finishes with
\texttt{field\_simp}. The first draft included a trailing \texttt{ring}
out of habit; \texttt{field\_simp} had already closed the goal,
producing the error ``\texttt{No goals to be solved}.''  Fix: drop the
\texttt{ring}. The lesson is the laplace \texttt{CLAUDE.md} entry on
``\texttt{field\_simp} sometimes finishes a proof, sometimes leaves a
ring residue'' — knowing which it does for a given goal is a matter
of experience, but the fix when it surprises you is one line.

\paragraph{Extracting $Z_k(t) = \mathrm{CONST}_k / t^\alpha$ from
$t^\alpha \cdot Z_k(t) = \mathrm{CONST}_k$.} The naive
\texttt{field\_simp at this; linarith} did not close the goal
\texttt{... = CONST / t\^{}α} because \texttt{linarith} could not unify
the resulting forms. Cleaner: \texttt{rw [eq\_div\_iff
ht\_rpow\_ne, mul\_comm]; exact hZt}. The \texttt{eq\_div\_iff} lemma
is the canonical Mathlib idiom for ``\texttt{A = B / C} iff
\texttt{A * C = B}''; combined with a \texttt{mul\_comm} to align with
the existing $t^\alpha \cdot Z_k(t)$ form, the proof is two rewrites.

\paragraph{\texttt{Filter.Eventually} has no \texttt{.symm}.} The
\texttt{Tendsto.congr'} step at the end of the unconditional theorem
takes an \texttt{EventuallyEq} and a \texttt{Tendsto}; my
\texttt{heq\_ratio} was stated in the direction
``\texttt{LHS = RHS}'', but \texttt{Tendsto.congr'} wanted the symmetric
direction. The reflex \texttt{heq\_ratio.symm} produces a type error
because \texttt{Filter.Eventually} doesn't have a field projection
\texttt{.symm} — \texttt{EventuallyEq} is a \texttt{def}, not a
structure. The canonical idiom is
\texttt{(heq\_ratio.mono fun \_ h => h.symm)}, which lifts the per-element
symmetry through the filter. Worth promoting to the laplace
\texttt{CLAUDE.md}.

\paragraph{Deliberation flag.} GPT-5.5~Pro's deliberation correctly
identified the structural risk in the unconditional form: \emph{``This
does \textbf{not} mean $J_k \sim 0$ in Mathlib's \texttt{IsEquivalent}
sense. The correct zero-case translation is $J_k = o[atTop]\,Z_k$.''}
That zero-case corollary is naturally a $\sim$5-line follow-up if any
downstream consumer needs it; deferred from this tide. The voting
converged on the first round.

\section{What was Mathlib, what was new}

\paragraph{Mathlib.} Standard tendsto chasing on the constant: the
two \texttt{Tendsto} combinators \texttt{div\_const} and
\texttt{congr'}, the rpow / Gamma / div\_rpow positivity lemmas, the
\texttt{eq\_div\_iff} idiom, and the \texttt{IsEquivalent}-to-ratio
bridge.\footnote{Specifically
\texttt{Real.rpow\_pos\_of\_pos}, \texttt{Real.Gamma\_pos\_of\_pos},
\texttt{Real.div\_rpow}, \texttt{eq\_div\_iff}, and
\texttt{Asymptotics.isEquivalent\_iff\_tendsto\_one}.}

\paragraph{New.} Two theorems and zero infrastructure. Everything
else --- the J-function asymptotic itself, the closed-form partition,
the positivity facts --- was already present from N1 and the broader
seabed.

The tide stays well-bounded in scope. The 200-line size is structural
overhead (positivity, constant-massaging, \texttt{Pi.div} cleanup) rather
than proof difficulty.

\section{Lessons}

\begin{enumerate}[leftmargin=*]

\item \textbf{Closed-form denominators turn ``asymptotic-equivalence''
into ``\texttt{Tendsto.div\_const}''.} When the denominator has an exact
closed form (as opposed to merely an asymptotic), the ratio of two
$t$-dependent quantities becomes a one-step \texttt{Tendsto} on the
numerator divided by a constant. This is much cleaner than chasing
two simultaneous asymptotics. Worth flagging for future tides where
the denominator is a partition function: don't reach for \texttt{Tendsto.div}
when \texttt{Tendsto.div\_const} suffices.

\item \textbf{The \texttt{IsEquivalent} idiom requires eventual
non-vanishing of the target.} Mathlib's bridge to ratio convergence is
\texttt{isEquivalent\_iff\_tendsto\_one}, which assumes
$g$ is eventually nonzero. For paper-style statements
$f \sim c \cdot g$ with $c$ a parameter that might be zero, the
\texttt{IsEquivalent} form requires $c \neq 0$ as a hypothesis;
the unconditional ratio form does not. So the natural
public API is the pair: unconditional ratio + IsEquivalent corollary
under non-vanishing of the leading coefficient. Both are useful for
different downstream contexts.

\item \textbf{\texttt{Filter.Eventually} symmetry needs
\texttt{.mono}, not \texttt{.symm}.} The reflex
\texttt{heq.symm} fails on \texttt{EventuallyEq} (it's a \texttt{def},
not a structure). The canonical idiom is
\texttt{(heq.mono fun \_ h => h.symm)} — lift the per-element symmetry
through the filter. Worth promoting to the laplace \texttt{CLAUDE.md}
under the existing ``\texttt{Pi.add} vs single-lambda'' entry; both
are filter-level vs. point-level mismatches.

\end{enumerate}

\section{Follow-ups}

\begin{itemize}[leftmargin=*]

\item \textbf{The $\pi(0) = 0$ corollary $J_k =o[atTop] Z_k$.} GPT
flagged this as the correct \emph{little-o} form for the zero-case.
$\sim$5 lines on top of the unconditional ratio theorem
(\texttt{Tendsto.div\_const} of $J_k/Z_k \to 0$ via \texttt{IsLittleO}'s
\texttt{Tendsto}-characterisation). Defer until a downstream consumer
asks; cheap to add then.

\item \textbf{The $k=2$ and $k=3$ specialisations.} The grammar
paper, the primer, and the L1 retrospective all reference the
quartic and sextic $\sim$-forms specifically. With the parametric
form in hand these are zero-line corollaries; named wrappers under
\texttt{quartic\_jfunction\_isEquivalent} and the obvious sextic
analogue would make the \texttt{\textbackslash leanref} story for those
papers clean. $\sim$10 lines each. Best landed once the relevant
paper-side has TeX statements to tag.

\item \textbf{Hypothesis weakening (L3 in the v3 survey).} The
present tide inherits N1's \texttt{Continuous prior + globally
bounded} hypothesis class. The minimal hypothesis class for the
underlying DCT step is
\texttt{AEMeasurable + ContinuousAt 0 + global bound}, per the L3
follow-up. Refactoring the parametric $J_k$ asymptotic to the weaker
class would automatically propagate to the
\texttt{\_isEquivalent} wrapper here. Bigger tide.

\item \textbf{Promote the partition closed form to a named
\texttt{kth\_partition\_smul\_pow} or analogous.} The rearrangement
$t^\alpha \cdot Z_k(t) = \mathrm{CONST}_k$ is a useful canonical form
for any tide that wants to reason about ratios of $J$-like
integrals. Pulling it out as a named one-line corollary of
\texttt{kth\_partition} would shave $\sim$15 lines off the present
file's \texttt{set} + algebra block, and become reusable. Cheap;
defer until a second consumer materialises.

\end{itemize}

\end{document}
